\section{Cycle classes in families}\label{sec:families}

Let $Y\to B$ be a family of almost-complex six-manifolds. The base $B$
is a finite simplicial complex. On each closed simplex the family is
a smooth submersion with a smooth vertical almost-complex structure,
in charts extending across the faces. The restrictions agree on common
faces. Open and empty targets are allowed. Write $Z(Y/B)$ for the space of
compact effective pseudoholomorphic one-cycles in the fibers, including
the empty cycle. Its topology records the base parameter, the Hausdorff
limit of the supports and weak convergence of their integration currents
\cite[Definition~2.6]{DIW}. Empty support is a separate open and closed
locus. A morphism of families is an inclusion of simplicial bases and
a fiberwise almost-complex open embedding into the restricted target,
smooth on every closed simplex.

For a locally compact Hausdorff space $\pi:Z\to B$ put
\[
             V(Z/B)=H_c^0(Z;\pi^*\omega_B),
\]
where $\omega_B$ is the rational dualizing complex. Cohomology here
and below is sheaf cohomology. In cellular form the summand over
an oriented $p$-simplex $\sigma$ consists of the compactly supported
cochains of $Z|_\sigma$ shifted by $p$. The second differential takes
signed restrictions to the faces. This follows by resolving $\omega_B$
by the constant sheaves on the closed simplices with the indicated shifts
\cite[Definitions~2.2 and~3.1, Lemma~3.4]{Pardon}.

\subsection{Topology and operations}
Over a small base chart we can identify a neighborhood of any compact
support with an open part of a product. Parallel transport on one
closed simplex gives such a trivialization after shrinking about the
compact set. A cutoff Hermitian form at the central parameter is
uniformly positive on nearby complex lines, so convergence of currents
bounds their masses in the chart. Consequently changes of trivialization
preserve convergence. Requiring the supports to lie in a relatively
compact neighborhood and imposing a closed bound on mass gives a
compact containing set by \cite[Theorem~2.11]{DIW}. For this application
one may extend the structures by cutoff to the fixed central fiber
on a slightly larger neighborhood. Strict inequalities give a
neighborhood of the prescribed cycle. The finite gluing along closed
face spaces gives a locally compact Hausdorff space over $B$.

On a complex analytic family the topology on each simplex agrees with
Barlet's topology on compact cycles. Apply the support and volume
criterion of \cite[\S1]{Barlet} in the complexification of the
simplex. Only the vertical $(1,1)$ part of a test two-form pairs with
a complex curve. Differences of positive Hermitian forms suffice for
these tests on a compact set. A positive form supported near a given
point of the limiting curve detects lower continuity of supports.
Conversely, compact cutoffs and uniform approximation on a set of
bounded mass give Barlet's continuous volume tests.

Let $S$ be the rational vector-space colimit of $V(Z(Y/B)/B)$ over
these families, identifying subdivisions. A closed base inclusion
includes the corresponding face summands. An open embedding of targets
extends by zero on the open subset of cycles contained in the source.
These two operations commute by restriction of extension by zero.
The external product, with disjoint union of the two targets over the
product base, defines a commutative multiplication on $S$. The empty
target over a point gives the unit. These are the smooth counterparts
of the operations of \cite[\S3.1]{Pardon}. Extension by zero from the
open semi-Fano cycle locus, followed by passage to the underlying
almost-complex family, thus gives a ring homomorphism
\[
                          \Phi:H\longrightarrow S .
\]
On the summand defining $H$ the relative cohomological degree is zero.

If the base is a point, the characteristic function of a compact open
subset of its cycle space gives a class of $S$. Such a class remembers
the ambient almost-complex manifold. An insertion on a family over
$B\times[0,1]$, with $B$ a closed oriented manifold of dimension $n$,
can also be represented by a compactly supported $n$-cocycle.
Restrict it to each top simplex and sum with the product orientations.
In the shifted complex this is a cochain of degree $-1$. Interior
faces cancel in its differential, leaving the difference of the two
endpoint insertions. We use this relation both for characteristic
functions over an interval and for Thom classes over a finite Borel
construction.

For the addition operation, retain all relative degrees in $V$ and
write $S^*$ for the resulting graded group. Proper pullback by addition
over the same base followed by pushforward along the diagonal of bases
takes values in $S^*\otimes S^*$, as in
\cite[Definitions~2.3 and~3.9]{Pardon}. Here the external product is
an isomorphism for the bivariant groups over two independent bases.
Indeed, for the compactly supported cohomology of two locally compact
Hausdorff spaces use their one-point compactifications relative to
the added points. On the compact product finite rectangular covers
are cofinal. The nerve of a product cover maps to the product of the
two nerves by taking marginal barycentric coordinates, with a section
given by their products. Linear interpolation gives a homotopy on
each product of simplices, also on the relative subspaces at infinity.
Finite-polyhedron Künneth over $\QQ$ therefore passes to relative
Čech cohomology, hence to compactly supported sheaf cohomology.
Filtering the two finite face complexes by dimension gives the
bivariant assertion, since the product is an isomorphism on the first
pages. Use product subdivisions containing the diagonal for the
base pushforward. Additivity of support of effective cycles makes
this construction commute with open extension, and restriction to
faces commutes with both operations. The same argument for
multiplication by $d$ gives $\rho_d$.
